\section{Comparação A/B com estudantes}
\label{sec:avaliacao-estudantes}

Esta avaliação foi feita com uma turma de estudantes, que comparou duas versões do mesmo programa: a Versão A, na sintaxe padrão, e a Versão B, personalizada em português. O formulário também traz afirmações sobre a plataforma e a demonstração.

\subsection{Formulário e procedimento}
\label{subsec:formulario-estudantes}

O formulário dos estudantes não constava da versão anterior do protocolo. Ele foi preparado para uma sessão de demonstração com uma turma do 1º ano do ensino médio técnico em Informática da UFCI-MG, Campus Leopoldina, realizada em 3 de novembro de 2026. Os autores demonstraram a plataforma em sala, os alunos também a usaram e responderam ao formulário em seguida.

O formulário é curto. Depois do termo de assentimento, pergunta a faixa de idade, o curso, o ano e a experiência anterior com programação. Seguem seis afirmações em escala de concordância de cinco pontos: se o aluno entendeu para que serve a plataforma, se um programa escrito com palavras em português é mais fácil de entender do que um escrito com palavras em inglês, se poder escolher as palavras da linguagem ajudaria quem está começando, se ver a separação do programa em \textit{tokens} ajuda a entender como o computador lê o código, se gostaria de usar a plataforma nas aulas e se a demonstração aumentou seu interesse em aprender programação. Uma questão mostra duas versões de um mesmo programa, que lê um número e escreve se ele é par ou ímpar. A versão A usa a sintaxe padrão, em inglês, e a versão B usa uma sintaxe personalizada em português. O aluno indica qual das duas acha mais fácil de entender, ou se as duas são igualmente fáceis ou igualmente difíceis. A medida obtida com essa pergunta é a preferência declarada entre as versões, isto é, o número de alunos que escolheu cada uma das quatro alternativas; ela registra o que o aluno diz achar mais fácil e não mede compreensão. O formulário termina com um campo de comentário livre.

Dois desses itens tocam o tema deste trabalho. O item sobre os \textit{tokens} leva ao público-alvo uma pergunta sobre a exibição de uma fase intermediária da tradução, que o formulário dos especialistas só alcança pela opinião de especialistas. A questão das duas versões depende de uma propriedade que o núcleo garante, a de que vocabulários diferentes podem descrever o mesmo programa. O que ela mede, porém, é qual das duas escritas o aluno declara entender com mais facilidade.

Os limites desse instrumento são maiores que os do formulário dos especialistas. A turma é uma só e foi escolhida por conveniência. A opinião foi colhida logo depois de uma demonstração feita pelos próprios autores e de um período de uso da plataforma, situação que favorece o efeito de novidade e a tendência a concordar. O formulário não mede desempenho. Ele serve como indício de aceitação e de preferência declarada.

\subsection{Resultados da sessão}
\label{sec:resultado-estudantes}

Esta subseção trata das afirmações sobre a demonstração (questões \ref{qF:demo-ini} a \ref{qF:demo-fim}) e da pergunta de comparação entre as versões (questão \ref{qF:comparacao}) do Apêndice~\ref{ap:formulario-estudantes}. O formulário descrito na Subseção~\ref{subsec:formulario-estudantes} foi aplicado em 3 de novembro de 2026 a uma turma do 1º ano do ensino médio técnico em Informática da UFCI-MG, Campus Leopoldina, depois de uma demonstração da plataforma feita pelos autores e de um período em que os alunos também a usaram. Houve dezoito respostas, todas com o termo de assentimento marcado. Nove alunos tinham até 15 anos e nove tinham 16 ou 17. Nove disseram já ter programado um pouco, oito disseram programar com frequência e um disse nunca ter programado. Quase ninguém, portanto, era iniciante absoluto.

A Tabela~\ref{tab:resultado-estudantes} traz as seis afirmações. Todas tiveram média entre 4,00 e 4,33. A afirmação de que um programa em português é mais fácil de entender do que um em inglês teve a menor média e o maior número de neutros: seis dos dezoito não tomaram posição sobre o item mais próximo da premissa da plataforma.

\begin{table}[htbp]
\centering
\caption{Afirmações do formulário da sessão com estudantes (escala de 1 a 5; dezoito respostas).}
\label{tab:resultado-estudantes}
\footnotesize
\renewcommand{\arraystretch}{1.2}
\begin{tabular}{|p{6.3cm}|c|c|c|c|c|}
\hline
\textbf{Afirmação} & \textbf{Média} & \textbf{Mediana} & \textbf{Concordam} & \textbf{Neutros} & \textbf{Discordam} \\\hline
Entendi para que serve a plataforma & 4,33 & 4,5 & 17 & 0 & 1 \\\hline
Um programa escrito com palavras em português é mais fácil de entender do que um escrito com palavras em inglês & 4,00 & 4 & 11 & 6 & 1 \\\hline
Poder escolher as palavras da linguagem ajudaria quem está começando a programar & 4,17 & 4 & 14 & 3 & 1 \\\hline
Ver como a plataforma separa o programa em partes (os \textit{tokens}) ajuda a entender como o computador lê o código & 4,17 & 4 & 15 & 1 & 2 \\\hline
Eu gostaria de usar essa plataforma nas aulas de programação & 4,17 & 4 & 15 & 3 & 0 \\\hline
A demonstração aumentou meu interesse em aprender programação & 4,17 & 4 & 16 & 0 & 2 \\\hline
\end{tabular}
\legend{Fonte: dados da pesquisa (2026).}
\end{table}

Para este trabalho, o item central é o dos \textit{tokens}, porque leva ao público-alvo uma pergunta sobre a exibição de uma fase da tradução. Quinze dos dezoito alunos concordaram que ver a separação do programa em \textit{tokens} ajuda a entender como o computador lê o código, um ficou no ponto neutro e dois discordaram. Os oito que programam com frequência concordaram todos. Entre os nove que já tinham programado um pouco, sete concordaram, um ficou no ponto neutro e um discordou. O aluno que nunca tinha programado também discordou. No painel, a mesma fase teve média 4,50 como recurso e 4,29 como valor didático, mas só dois especialistas declararam ter visto a tabela. Os estudantes a viram na demonstração, e a concordância deles é o dado mais direto que o trabalho tem sobre essa fase. Ainda assim, o item registra concordância com uma frase e não mede compreensão.

O formulário também mostrou duas versões de um programa que lê um número e escreve se ele é par ou ímpar: a Versão A, na sintaxe padrão em inglês, e a Versão B, personalizada em português. A Tabela~\ref{tab:resultado-versoes} separa as respostas pela experiência declarada.

\begin{table}[htbp]
\centering
\caption{Versão do programa considerada mais fácil de entender, por experiência declarada com programação.}
\label{tab:resultado-versoes}
\footnotesize
\renewcommand{\arraystretch}{1.2}
\begin{tabular}{|p{4.3cm}|c|c|c|c|c|}
\hline
\textbf{Experiência} & \textbf{Versão B} & \textbf{Igualmente fáceis} & \textbf{Versão A} & \textbf{Igualmente difíceis} & \textbf{Total} \\\hline
Programa com frequência & 6 & 2 & 0 & 0 & 8 \\\hline
Já programou um pouco & 5 & 2 & 2 & 0 & 9 \\\hline
Nunca programou & 0 & 0 & 0 & 1 & 1 \\\hline
\textbf{Total} & 11 & 4 & 2 & 1 & 18 \\\hline
\end{tabular}
\legend{Fonte: dados da pesquisa (2026).}
\end{table}

Onze alunos acharam a Versão B mais fácil de entender, quatro acharam as duas igualmente fáceis, dois preferiram a Versão A e um achou as duas igualmente difíceis. A questão depende de uma propriedade que o núcleo garante: as duas versões são o mesmo programa. As palavras da Versão B são convertidas nos mesmos \textit{tokens} que as da Versão A, e a partir daí a análise sintática, o código intermediário e a execução não distinguem as duas. O formulário afirmava ao aluno que as versões fazem a mesma coisa, e o que ele mediu foi qual das duas escritas o aluno declara entender com mais facilidade. Não mediu se o aluno reconheceria a equivalência sozinho.

A preferência pela versão em português apareceu nos dois grupos com experiência. Os dois alunos que preferiram a versão em inglês estão entre os que já tinham programado um pouco, e um deles é o que discordou do item dos \textit{tokens}. O único aluno que nunca tinha programado, justamente o perfil para o qual a plataforma foi pensada, achou as duas versões igualmente difíceis e discordou de três das seis afirmações, entre elas a de ter entendido para que serve a plataforma e a dos \textit{tokens}. Uma única resposta não sustenta conclusão, mas ela vai em sentido contrário ao das demais e fica registrada.

O campo de comentário livre recebeu duas respostas: uma elogiou a plataforma e a outra pediu o acréscimo de outra linguagem. O formulário não tem item de autorização de citação, e por isso os comentários não são transcritos.

Os limites dessa sessão estão na Subseção~\ref{subsec:formulario-estudantes}. Foi uma única turma, de conveniência, que respondeu logo depois de uma demonstração feita pelos próprios autores e de um período de uso da plataforma, e quase todos já programavam. Nenhuma medida de desempenho foi colhida. As respostas servem como indício de aceitação e de preferência declarada pela versão em português. Não mostram redução da barreira sintática nem aprendizagem. O registro do assentimento e a falta de autorização dos responsáveis estão descritos na Seção~\ref{sec:protocolo-avaliacao} e retomados nas limitações do Capítulo~\ref{cap:consideracoes}.

